% ==============================================================================
% ==============================================================================
% Documento de Especificación del Proyecto - Sistema de Ventas B2B
% ==============================================================================
% ==============================================================================

\documentclass[12pt,a4paper,oneside]{book}

% -----------------------------------------------------------------------------
% Paquetes y configuración
% -----------------------------------------------------------------------------
\usepackage[spanish,es-tabla]{babel}       % Idioma español
\usepackage[utf8]{inputenc}                  % Codificación UTF-8
\usepackage[T1]{fontenc}                     % Codificación de fuentes
\usepackage{graphicx}                        % Gráficos e imágenes
\usepackage{hyperref}                        % Hipervínculos
\usepackage{geometry}                        % Configuración de página
\usepackage{titlesec}                        % Formato de títulos
\usepackage{fancyhdr}                        % Encabezados y pies
\usepackage{booktabs}                        % Tablas profesionales
\usepackage{listings}                        % Código fuente
\usepackage{xcolor}                          % Colores
\usepackage{amsmath,amssymb}                 % Matemáticas
\usepackage{parskip}                         % Espaciado entre párrafos
\usepackage{enumitem}                        % Listas personalizadas
\usepackage{pdflscape}                       % Páginas horizontales
\usepackage{longtable}                       % Tablas largas
\usepackage{rotating}                        % Rotación de elementos

% Configuración de hipervínculos
\hypersetup{
    colorlinks=true,
    linkcolor=blue,
    filecolor=magenta,
    urlcolor=cyan,
    pdftitle={Sistema de Ventas B2B - Especificación Técnica},
    pdfauthor={Equipo de Desarrollo},
    pdfsubject={Proyecto de Plataforma de Ventas B2B},
    pdfkeywords={B2B, Spring Boot, React, E-commerce, JWT, 2FA}
}

% Configuración de página
\geometry{
    left=25mm,
    right=25mm,
    top=30mm,
    bottom=25mm,
    headheight=15pt,
    footskip=15pt
}

% Colores personalizados
\definecolor{primary}{RGB}{0,82,147}
\definecolor{secondary}{RGB}{52,152,219}
\definecolor{accent}{RGB}{231,76,60}
\definecolor{success}{RGB}{39,174,96}
\definecolor{dark}{RGB}{44,62,80}
\definecolor{light}{RGB={236,240,241}
\definecolor{codebg}{RGB}{245,245,245}
\definecolor{codekeyword}{RGB}{0,0,255}
\definecolor{codestring}{RGB}{163,21,21}

% Configuración de listados de código
\lstset{
    backgroundcolor=\color{codebg},
    basicstyle=\ttfamily\small,
    keywordstyle=\color{codekeyword}\bfseries,
    stringstyle=\color{codestring},
    commentstyle=\color{gray}\itshape,
    numbers=left,
    numberstyle=\tiny\color{gray},
    stepnumber=1,
    numbersep=5pt,
    tabsize=2,
    breaklines=true,
    showstringspaces=false,
    frame=single,
    framerule=0.5pt
}

% Configuración de títulos de secciones
\titleformat{\chapter}[display]
    {\normalfont\Huge\bfseries\color{primary}}
    {\chaptertitlename\ \thechapter}{20pt}{\Huge}
\titleformat{\section}
    {\normalfont\Large\bfseries\color{dark}}
    {\thesection}{1em}{}
\titleformat{\subsection}
    {\normalfont\large\bfseries\color{dark}}
    {\thesubsection}{1em}{}
\titleformat{\subsubsection}
    {\normalfont\normalsize\bfseries\color{dark}}
    {\thesubsubsection}{1em}{}

% Encabezados y pies de página
\pagestyle{fancy}
\fancyhf{}
\fancyhead[L]{\nouppercase{\leftmark}}
\fancyhead[R]{\thepage}
\fancyfoot[C]{\textbf{Sistema de Ventas B2B}}
\fancyfoot[R]{\thepage}
\fancyfoot[L]{\href{https://github.com}{Repositorio del Proyecto}}

% Comandos personalizados
\newcommand{\chaptercolor}[1]{\textcolor{primary}{#1}}
\newcommand{\alert}[1]{\textcolor{accent}{\textbf{#1}}}
\newcommand{\success}[1]{\textcolor{success}{\textbf{#1}}}
\newcommand{\info}[1]{\textcolor{secondary}{\textbf{#1}}}
\newcommand{\HRule}{\rule{\linewidth}{0.5mm}}

% Configuración de enumerate
\setlist[enumerate]{itemsep=0mm, partopsep=0mm}
\setlist[itemize]{itemsep=0mm, partopsep=0mm}

% ==============================================================================
% INICIO DEL DOCUMENTO
% ==============================================================================

\begin{document}

% -----------------------------------------------------------------------------
% PORTADA
% -----------------------------------------------------------------------------
\begin{titlepage}
    \centering
    \vspace*{2cm}
    
    % Logo o imagen
    \includegraphics[width=0.3\textwidth]{}\\[2cm]
    \textbf{\LARGE\color{primary}Sistema de Ventas B2B}\\[0.5cm]
    \textbf{\large con Pasarela de Pago}\\[2cm]
    
    \HRule\\[0.5cm]
    \textbf{\huge Documento de Especificación Técnica}\\[0.5cm]
    \HRule\\[2cm]
    
    \begin{table}[ht]
        \centering
        \begin{tabular}{p{3.5cm}p{10cm}}
            \textbf{Backend:} & Java + Spring Boot \\
            \textbf{Frontend:} & React + Vite \\
            \textbf{Base de datos:} & MySQL/MariaDB o PostgreSQL \\
            \textbf{Versión:} & 1.0 \\
            \textbf{Fecha:} & \today
        \end{tabular}
    \end{table}\\[2cm]
    
    \vspace*{\fill}
    \textbf{\textit{Proyecto de Desarrollo - 2026}}\\[0.5cm]
\end{titlepage}

% -----------------------------------------------------------------------------
% PÁGINA DE LICENCIA Y INFO
% -----------------------------------------------------------------------------
\clearpage
\thispagestyle{empty}
\vspace*{5cm}
\begin{center}
    \textbf{\Large Acerca de este Documento}\\[1cm]
    
    \begin{minipage}{0.8\textwidth}
        Este documento contiene la especificación técnica completa para el desarrollo 
        del Sistema de Ventas B2B. Incluye los objetivos del proyecto, alcance funcional, 
        arquitectura, modelo de datos y flujos del sistema.
    \end{minipage}\\[2cm]
    
    \HRule\\[1cm]
    \begin{tabular}{p{5cm}p{8cm}}
        \textbf{Estado:} & En Desarrollo \\
        \textbf{Tipo:} & Especificación Técnica \\
        \textbf{Clasificación:} & Público
    \end{tabular}
\end{center}

% -----------------------------------------------------------------------------
% TABLA DE CONTENIDOS
% -----------------------------------------------------------------------------
\clearpage
\tableofcontents
\clearpage

% ==============================================================================
% CAPÍTULO 1: OBJETIVOS DEL PROYECTO
% ==============================================================================
\chapter{Objetivos del Proyecto}
\label{cap:objetivos}

\section{Propósito del Sistema}
El presente documento tiene como propósito establecer las especificaciones técnicas 
y funcionales para el desarrollo de una plataforma web de ventas B2B (Business-to-Business) 
que permita a empresas registradas gestionar sus pedidos de manera eficiente y segura.

\section{Objetivos Estratégicos}

\subsection{Objetivo General}
Desarrollar una plataforma web de ventas B2B que permita:
\begin{itemize}
    \item El registro y gestión de empresas en el sistema.
    \item La administración de múltiples usuarios por empresa.
    \item La autenticación segura mediante JWT con soporte para 2FA.
    \item El CRUD completo de productos y categorías.
    \item La gestión de pedidos y-carrito de compras.
    \item El procesamiento de pagos mediante pasarela externa.
    \item La gestión de contenido (CMS) desde el panel de administración.
    \item El registro de auditoría y seguridad.
\end{itemize}

\section{Alcance del Proyecto}

\subsection{Incluido en el Sistema}
\begin{itemize}
    \item Registro de empresas con NIT y razón social.
    \item Múltiples usuarios por empresa con roles diferenciados.
    \item Autenticación JWT con refresh tokens.
    \item Autenticación de dos factores (2FA) mediante TOTP.
    \item Recuperación y cambio de contraseña.
    \item Cerrar sesión segura con invalidación de tokens.
    \item CRUD de productos y categorías (gestión desde admin).
    \item Subida de imágenes múltiples por producto.
    \item Carrito de compras persistente.
    \item Lista de deseos (favoritos).
    \item Cotizaciones B2B.
    \item Pedidos recurrentes (compras programadas).
    \item Pasarela de pago (Stripe/MercadoPago).
    \item Confirmación automática de compras.
    \item Gestión de contenido (CMS).
    \item Registro de auditoría.
\end{itemize}

\subsection{No Incluido en el Sistema}
\begin{alert}{IMPORTANTE:} El sistema \textbf{NO} maneja suscripciones ni cobros mensuales. 
Solo procesamiento de compras directas por pedido.

\section{Resultados Esperados}
\begin{enumerate}
    \item Plataforma web funcional para ventas B2B.
    \item Sistema de autenticación seguro con JWT y 2FA.
    \item Panel de administración completo.
    \item Integración con pasarela de pago.
    \item Documentación técnica actualizada.
\end{enumerate}

% ==============================================================================
% CAPÍTULO 2: ALCANCE FUNCIONAL
% ==============================================================================
\chapter{Alcance Funcional}
\label{cap:alcance}

\section{Autenticación y Seguridad}
\label{sec:autenticacion}

\subsection{Registro de Usuarios}
Sistema de registro que permite:
\begin{itemize}
    \item Registro de empresas con datos completos (NIT, razón social).
    \item Registro de usuarios asociados a empresas.
    \item Validación de datos única (email, NIT, username).
\end{itemize}

\subsection{Inicio de Sesión JWT}
\begin{itemize}
    \item Generación de access token (duración: 15-30 minutos).
    \item Generación de refresh token (duración: 7 días).
    \item Renovación de sesión automática mediante refresh token.
\end{itemize}

\subsection{Autenticación de Dos Factores (2FA)}
\begin{info}{CARACTERÍSTICA CLAVE:} Implementación de TOTP (Time-based One-Time Password) 
compatible con Google Authenticator y aplicaciones similares.

\begin{itemize}
    \item Habilitación opcional desde el perfil del usuario.
    \item Generación de secreto TOTP único por usuario.
    \item Visualización de código QR para escaneo.
    \item Verificación de código inicial para activar 2FA.
    \item Solicitud de código en cada inicio de sesión si está habilitado.
\end{itemize}

\subsection{Recuperación de Contraseña}
\begin{enumerate}
    \item Solicitud mediante email.
    \item Generación de token de recuperación (expiración: 1 hora).
    \item Envío de email con enlace de recuperación.
    \item Validación de token y actualización de contraseña.
    \item Invalidación de tokens anteriores.
\end{enumerate}

\subsection{Cambio de Contraseña}
\begin{itemize}
    \item Cambio desde el perfil de usuario.
    \item Validación de contraseña actual.
    \item Invalidación de todos los tokens activos (logout de todos los dispositivos).
    \item Envío de notificación por email.
\end{itemize}

\subsection{Cerrar Sesión (Logout)}
\begin{enumerate}
    \item Solicitud de cierre de sesión.
    \item Marcado de refresh token como revocado.
    \item Eliminación de access token del cliente.
    \item (Opcional) Cierre de sesiones en todos los dispositivos.
\end{enumerate}

\subsection{Roles y Permisos}
\begin{table}[ht]
    \centering
    \caption{Roles del sistema y permisos asociados}
    \begin{tabular}{@{}ll@{}}
        \toprule
        \textbf{Rol} & \textbf{Descripción} \\
        \midrule
        ADMIN & Administrador del sistema \\
        COMPANY\_ADMIN & Administrador de empresa \\
        BUYER & Comprador/Usuario de empresa \\
        CONTENT\_MANAGER & Gestor de contenido CMS \\
        \bottomrule
    \end{tabular}
\end{table}

\section{Gestión de Empresas}
\label{sec:empresas}

\subsection{Registro de Empresas}
\begin{itemize}
    \item NIT (Número de Identificación Tributaria).
    \item Razón social.
    \item Email de contacto.
    \item Teléfono.
    \item Dirección.
    \item Logo (opcional).
\end{itemize}

\subsection{Gestión de Empresa (Admin)}
\begin{itemize}
    \item CRUD completo de empresas.
    \item Activación/Desactivación de empresas.
    \item Visualización de empresas registradas.
    \item Gestión de usuarios por empresa.
\end{itemize}

\section{Ventas - Productos}
\label{sec:productos}

\subsection{Catálogo de Productos}
\begin{itemize}
    \item Nombre, descripción, precio base, stock.
    \item SKU único por producto.
    \item Peso y dimensiones.
    \item Estados: activo/inactivo.
\end{itemize}

\subsection{Categorías}
\begin{itemize}
    \item CRUD completo de categorías.
    \item Soporte para subcategorías (jerarquía).
    \item Imagen por categoría.
    \item Orden de visualización.
\end{itemize}

\subsection{Imágenes de Productos}
\begin{itemize}
    \item Múltiples imágenes por producto.
    \item Imagen principal (primary).
    \item Orden de visualización (sort\_order).
\end{itemize}

\subsection{Atributos Variables}
\begin{itemize}
    \item Tallas (S, M, L, XL, etc.).
    \item Colores.
    \item Otras variantes de producto.
\end{itemize}

\section{Carrito y Pedidos}
\label{sec:ventas}

\subsection{Carrito de Compras}
\begin{itemize}
    \item \success{CARRITO PERSISTENTE}: Se guarda aunque el usuario cierre sesión.
    \item Guardar carrito para después.
    \item Actualización de cantidades.
    \item Eliminación de items.
\end{itemize}

\subsection{Lista de Deseos (Wishlist)}
\begin{itemize}
    \item Agregar productos como favoritos.
    \item Visualización de productos deseados.
    \item Mover productos al carrito.
\end{itemize}

\subsection{Cotizaciones (B2B)}
\begin{itemize}
    \item Solicitar cotización sin comprar inmediatamente.
    \item Precios especiales por empresa.
    \item Validez de cotización configurable.
    \item Estados: PENDING, APPROVED, REJECTED, EXPIRED.
\end{itemize}

\subsection{Pedidos}
\begin{itemize}
    \item Creación de pedido desde carrito.
    \item Estados del pedido: CART, PENDING\_PAYMENT, PAID, SHIPPED, DELIVERED, CANCELLED.
    \item Dirección de envío.
    \item Notas adicionales.
    \item Historial de compras.
\end{itemize}

\section{Descuentos y Precios}
\label{sec:descuentos}

\subsection{Descuentos por Cantidad (Price Tiers)}
\begin{itemize}
    \item Precios por volumen/configuración de rangos.
    \item Ejemplo: 1-10 unidades = \$10, 11-50 = \$9, etc.
    \item Almacenamiento en formato JSON.
\end{itemize}

\subsection{Descuentos por Empresa (B2B)}
\begin{itemize}
    \item Precios especiales definidos por empresa.
    \item Acuerdos comerciales personalizados.
\end{itemize}

\subsection{Cupones de Descuento}
\begin{itemize}
    \item Código de cupón.
    \item Porcentaje o monto fijo.
    \item Validez por fecha.
    \item Usos limitados.
\end{itemize}

\subsection{Límites de Crédito}
\begin{itemize}
    \item Comprar ahora, pagar después.
    \item Configuración de crédito por empresa.
    \item Control de deuda pendiente.
\end{itemize}

\section{Gestión de Contenido (CMS)}
\label{sec:cms}

Desde el panel de administración se puede gestionar:

\begin{itemize}
    \item \textbf{Página principal}: Hero, banners, productos destacados.
    \item \textbf{About Us / Nosotros}: Texto institucional.
    \item \textbf{Información de contacto}: Datos de contacto.
    \item \textbf{Políticas}: Términos y condiciones, privacidad.
    \item \textbf{Footer y encabezados}: Configuración visual.
    \item \textbf{Banners promocionales}: Imágenes rotativas.
    \item \textbf{Noticias / Blog}: Artículos informativos.
\end{itemize}

\section{Archivos y Multimedia}
\label{sec:archivos}

\begin{itemize}
    \item Subida de imágenes.
    \item Imágenes para productos.
    \item Avatar de usuario.
    \item Galerías multimedia.
    \item Almacenamiento en S3 o Cloudinary.
\end{itemize}

\section{Seguridad del Sistema}
\label{sec:seguridad}

\begin{itemize}
    \item JWT con access token y refresh token.
    \item Hash de contraseñas (bcrypt/Argon2id).
    \item 2FA (TOTP).
    \item Roles y permisos granulares.
    \item Auditoría de acciones (AuditLog).
    \item Rate limiting (límite de peticiones).
    \item Protección CSRF.
    \item HTTPS obligatorio.
\end{itemize}

% ==============================================================================
% CAPÍTULO 3: STACK TECNOLÓGICO
% ==============================================================================
\chapter{Stack Tecnológico}
\label{cap:stack}

\section{Backend}
\label{sec:backend}

\begin{table}[ht]
    \centering
    \caption{Tecnologías del Backend}
    \begin{tabular}{@{}llp{8cm}@{}}
        \toprule
        \textbf{Tecnología} & \textbf{Versión} & \textbf{Descripción} \\
        \midrule
        Java & 17 & Lenguaje de programación \\
        Spring Boot & 3.5.10 & Framework principal \\
        Spring Security & 6.x & Seguridad y autenticación \\
        JWT (jjwt) & 0.12.5 & Generación y validación de tokens \\
        JPA/Hibernate & - & ORM para acceso a datos \\
        MySQL/MariaDB & 8.x & Base de datos (via XAMPP) \\
        PostgreSQL & 15+ & Base de datos alternativa \\
        MercadoPago SDK & 2.1.1 & Pasarela de pago \\
        S3/Cloudinary & - & Almacenamiento de archivos \\
        dev.samstevens.totp & 1.7.1 & Librería TOTP para 2FA \\
        Lombok & - & Reducción de boilerplate \\
        Bucket4j & 8.10.1 & Rate limiting \\
        SpringDoc OpenAPI & 2.3.0 & Documentación Swagger UI \\
        \bottomrule
    \end{tabular}
\end{table}

\section{Frontend}
\label{sec:frontend}

\begin{table}[ht]
    \centering
    \caption{Tecnologías del Frontend}
    \begin{tabular}{@{}llp{8cm}@{}}
        \toprule
        \textbf{Tecnología} & \textbf{Versión} & \textbf{Descripción} \\
        \midrule
        React & 18+ & Biblioteca de UI \\
        Vite & 5+ & Build tool y servidor de desarrollo \\
        React Router & 6.x & Enrutamiento de páginas \\
        Axios & 1.x & Cliente HTTP \\
        Context API / Redux & - & Gestión de estado \\
        Tailwind CSS / MUI & - & Estilos y componentes \\
        react-qr-code & - & Generación de QR para 2FA \\
        \bottomrule
    \end{tabular}
\end{table}

\section{Infraestructura}
\label{sec:infraestructura}

\begin{itemize}
    \item \textbf{Docker}: Contenedores para despliegue.
    \item \textbf{Nginx}: Servidor web y proxy reverso.
    \item \textbf{HTTPS}: Seguridad en tránsito.
    \item \textbf{Email Service}: SendGrid, AWS SES o SMTP.
\end{itemize}

% ==============================================================================
% CAPÍTULO 4: ARQUITECTURA
% ==============================================================================
\chapter{Arquitectura del Sistema}
\label{cap:arquitectura}

\section{Arquitectura N-Capas por Módulo}
\label{sec:arquitectura_n_capas}

El proyecto sigue una arquitectura N-Capas por módulo, similar a una arquitectura de microservicios, 
pero monolítica. Cada módulo tiene su propia estructura de 4 capas:

\begin{figure}[ht]
    \centering
    \includegraphics[width=0.8\textwidth]{}
    \caption{Arquitectura N-Capas por Módulo}
    \label{fig:arquitectura_capas}
\end{figure}

\subsection{Estructura de Directorios}

\begin{lstlisting}[language=bash, numbers=none]
ventas-api/src/main/java/com/app/ventas_api/
|
+-- organizacion/              # MODULO: Empresas
|   +-- domain/               # Entidades, objetos de valor
|   +-- application/           # Casos de uso, servicios
|   +-- infrastructure/       # Repositorios, BD
|   +-- presentation/         # Controladores, DTOs
|
+-- productos/                 # MODULO: Productos
|   +-- domain/
|   +-- application/
|   +-- infrastructure/
|   +-- presentation/
|
+-- ventas/                   # MODULO: Pedidos y Pagos
|   +-- domain/
|   +-- application/
|   +-- infrastructure/
|   +-- presentation/
|
+-- seguridad/                # MODULO: Auth, JWT, Roles
|   +-- domain/              # User, Role, RefreshToken
|   +-- application/
|   +-- infrastructure/
|   +-- presentation/
|
+-- shared/                   # Comun
    +-- kernel/              # Utilidades globales
    +-- infrastructure/      # Config global
\end{lstlisting}

\subsection{Descripción de Capas}

\begin{description}
    \item[Domain] Entidades del negocio, objetos de valor, interfaces de repositorio.
    \item[Application] Casos de uso, servicios de aplicación, DTOs.
    \item[Infrastructure] Implementaciones de repositorios, configuración de BD.
    \item[Presentation] Controladores REST, DTOs de request/response.
\end{description}

\section{Modelo de Despliegue}
\label{sec:despliegue}

\begin{figure}[ht]
    \centering
    \begin{tikzpicture}[node distance=2cm, scale=0.8, transform shape]
        \node (Client) [rectangle, draw, fill=blue!20] {Navegador Cliente};
        \node (Nginx) [rectangle, draw, fill=green!20, right=of Client] {Nginx (Proxy)};
        \node (Backend) [rectangle, draw, fill=yellow!20, right=of Nginx] {Spring Boot API};
        \node (DB) [cylinder, draw, fill=orange!20, right=of Backend] {MySQL/PostgreSQL};
        \node (Storage) [rectangle, draw, fill=purple!20, below=of Backend] {S3/Cloudinary};
        \node (Payment) [rectangle, draw, fill=red!20, below=of Nginx] {Pasarela Pago};
        
        \draw [->] (Client) -- (Nginx);
        \draw [->] (Nginx) -- (Backend);
        \draw [->] (Backend) -- (DB);
        \draw [->] (Backend) -- (Storage);
        \draw [->] (Backend) -- (Payment);
    \end{tikzpicture}
    \caption{Arquitectura de Despliegue}
    \label{fig:despliegue}
\end{figure}

% ==============================================================================
% CAPÍTULO 5: MODELO ENTIDAD-RELACIÓN
% ==============================================================================
\chapter{Modelo Entidad-Relación (MER)}
\label{cap:mer}

\section{Diagrama ER Conceptual}
\label{sec:er_conceptual}

\begin{figure}[ht]
    \centering
    \begin{tikzpicture}[node distance=1.5cm, scale=0.7, transform shape]
        % Company
        \node (Company) [rectangle, draw, fill=blue!20] {Company};
        \node (User) [rectangle, draw, fill=blue!20, right=of Company] {User};
        \node (Role) [rectangle, draw, fill=blue!20, right=of User] {Role};
        
        % Products
        \node (Category) [rectangle, draw, fill=green!20, below=of Company] {Category};
        \node (Product) [rectangle, draw, fill=green!20, right=of Category] {Product};
        \node (ProductImage) [rectangle, draw, fill=green!20, right=of Product] {ProductImage};
        
        % Orders
        \node (Order) [rectangle, draw, fill=yellow!20, below=of Category] {Order};
        \node (OrderItem) [rectangle, draw, fill=yellow!20, right=of Order] {OrderItem};
        \node (Payment) [rectangle, draw, fill=yellow!20, right=of OrderItem] {Payment};
        
        % Relationships
        \draw [->] (Company) -- (User);
        \draw [->] (User) .. controls +(0,-1) and +(0,1) .. (Role);
        \draw [->] (Company) -- (Category);
        \draw [->] (Category) -- (Product);
        \draw [->] (Product) -- (ProductImage);
        \draw [->] (Company) -- (Order);
        \draw [->] (Order) -- (OrderItem);
        \draw [->] (OrderItem) -- (Product);
        \draw [->] (Order) -- (Payment);
    \end{tikzpicture}
    \caption{Diagrama Entidad-Relación Simplificado}
    \label{fig:er_simplificado}
\end{figure}

\section{Entidades del Sistema}
\label{sec:entidades}

A continuación se detallan todas las entidades del sistema con sus atributos y relaciones.

% -----------------------------------------------------------------------------
% COMPANY
% -----------------------------------------------------------------------------
\subsection{Company (Empresa)}
\label{sec:company}

\begin{table}[ht]
    \centering
    \caption{Tabla: Company}
    \begin{tabular}{@{}lll@{}}
        \toprule
        \textbf{Atributo} & \textbf{Tipo} & \textbf{Descripción} \\
        \midrule
        id & BIGINT (PK) & Identificador único \\
        nit & VARCHAR(20) & NIT de la empresa \\
        business\_name & VARCHAR(255) & Razón social \\
        email & VARCHAR(255) & Email de contacto \\
        phone & VARCHAR(20) & Teléfono \\
        address & VARCHAR(500) & Dirección \\
        logo\_url & VARCHAR(500) & URL del logo \\
        active & BOOLEAN & Estado activo/inactivo \\
        created\_at & DATETIME & Fecha de creación \\
        updated\_at & DATETIME & Fecha de actualización \\
        \bottomrule
    \end{tabular}
\end{table}

\textbf{Relaciones:}
\begin{itemize}
    \item 1:N con User (una empresa tiene muchos usuarios).
    \item 1:N con Order (una empresa tiene muchos pedidos).
\end{itemize}

% -----------------------------------------------------------------------------
% USER
% -----------------------------------------------------------------------------
\subsection{User (Usuario)}
\label{sec:user}

\begin{table}[ht]
    \centering
    \caption{Tabla: User}
    \begin{tabular}{@{}lll@{}}
        \toprule
        \textbf{Atributo} & \textbf{Tipo} & \textbf{Descripción} \\
        \midrule
        id & BIGINT (PK) & Identificador único \\
        company\_id & BIGINT (FK) & Empresa asociada \\
        username & VARCHAR(50) & Nombre de usuario \\
        email & VARCHAR(255) & Email único \\
        password\_hash & VARCHAR(255) & Hash de contraseña \\
        phone & VARCHAR(20) & Teléfono \\
        avatar\_url & VARCHAR(500) & URL del avatar \\
        two\_factor\_enabled & BOOLEAN & 2FA habilitado \\
        two\_factor\_secret & VARCHAR(255) & Secreto TOTP \\
        password\_reset\_token & VARCHAR(255) & Token recuperación \\
        password\_reset\_expires & DATETIME & Expiración token \\
        active & BOOLEAN & Estado activo/inactivo \\
        created\_at & DATETIME & Fecha de creación \\
        updated\_at & DATETIME & Fecha de actualización \\
        last\_login\_at & DATETIME & Último inicio de sesión \\
        \bottomrule
    \end{tabular}
\end{table}

\textbf{Relaciones:}
\begin{itemize}
    \item N:1 con Company.
    \item N:M con Role (vía UserRole).
    \item 1:N con Media.
    \item 1:N con AuditLog.
\end{itemize}

% -----------------------------------------------------------------------------
% ROLE
% -----------------------------------------------------------------------------
\subsection{Role (Rol)}
\label{sec:role}

\begin{table}[ht]
    \centering
    \caption{Tabla: Role}
    \begin{tabular}{@{}lll@{}}
        \toprule
        \textbf{Atributo} & \textbf{Tipo} & \textbf{Descripción} \\
        \midrule
        id & BIGINT (PK) & Identificador único \\
        name & VARCHAR(50) & Nombre del rol \\
        description & VARCHAR(255) & Descripción \\
        \bottomrule
    \end{tabular}
\end{table}

\textbf{Valores posibles:} ADMIN, COMPANY\_ADMIN, BUYER, CONTENT\_MANAGER

% -----------------------------------------------------------------------------
% USERROLE
% -----------------------------------------------------------------------------
\subsection{UserRole (Usuario-Rol)}
\label{sec:userrole}

\begin{table}[ht]
    \centering
    \caption{Tabla: UserRole (intermedia)}
    \begin{tabular}{@{}lll@{}}
        \toprule
        \textbf{Atributo} & \textbf{Tipo} & \textbf{Descripción} \\
        \midrule
        user\_id & BIGINT (FK) & Usuario \\
        role\_id & BIGINT (FK) & Rol \\
        \bottomrule
    \end{tabular}
\end{table}

% -----------------------------------------------------------------------------
% CATEGORY
% -----------------------------------------------------------------------------
\subsection{Category (Categoría)}
\label{sec:category}

\begin{table}[ht]
    \centering
    \caption{Tabla: Category}
    \begin{tabular}{@{}lll@{}}
        \toprule
        \textbf{Atributo} & \textbf{Tipo} & \textbf{Descripción} \\
        \midrule
        id & BIGINT (PK) & Identificador único \\
        name & VARCHAR(100) & Nombre de categoría \\
        description & TEXT & Descripción \\
        parent\_id & BIGINT (FK) & Categoría padre (subcategorías) \\
        image\_url & VARCHAR(500) & Imagen de categoría \\
        active & BOOLEAN & Estado activo/inactivo \\
        created\_at & DATETIME & Fecha de creación \\
        \bottomrule
    \end{tabular}
\end{table}

\textbf{Relaciones:}
\begin{itemize}
    \item 1:N con Product (una categoría tiene muchos productos).
    \item Autorreferencia para subcategorías.
\end{itemize}

% -----------------------------------------------------------------------------
% PRODUCT
% -----------------------------------------------------------------------------
\subsection{Product (Producto)}
\label{sec:product}

\begin{table}[ht]
    \centering
    \caption{Tabla: Product}
    \begin{tabular}{@{}lll@{}}
        \toprule
        \textbf{Atributo} & \textbf{Tipo} & \textbf{Descripción} \\
        \midrule
        id & BIGINT (PK) & Identificador único \\
        category\_id & BIGINT (FK) & Categoría asociada \\
        name & VARCHAR(255) & Nombre del producto \\
        description & TEXT & Descripción detallada \\
        base\_price & DECIMAL(10,2) & Precio base \\
        price\_tiers & JSON & Precios por cantidad \\
        stock & INT & Cantidad en inventario \\
        sku & VARCHAR(50) & Código único de producto \\
        weight & DECIMAL(10,2) & Peso (kg) \\
        dimensions & VARCHAR(50) & Dimensiones \\
        active & BOOLEAN & Estado activo/inactivo \\
        created\_at & DATETIME & Fecha de creación \\
        updated\_at & DATETIME & Fecha de actualización \\
        \bottomrule
    \end{tabular}
\end{table}

\textbf{Relaciones:}
\begin{itemize}
    \item N:1 con Category.
    \item 1:N con ProductImage.
    \item 1:N con OrderItem.
\end{itemize}

% -----------------------------------------------------------------------------
% PRODUCTIMAGE
% -----------------------------------------------------------------------------
\subsection{ProductImage (Imagen de Producto)}
\label{sec:productimage}

\begin{table}[ht]
    \centering
    \caption{Tabla: ProductImage}
    \begin{tabular}{@{}lll@{}}
        \toprule
        \textbf{Atributo} & \textbf{Tipo} & \textbf{Descripción} \\
        \midrule
        id & BIGINT (PK) & Identificador único \\
        product\_id & BIGINT (FK) & Producto asociado \\
        url & VARCHAR(500) & URL de la imagen \\
        is\_primary & BOOLEAN & Es imagen principal \\
        sort\_order & INT & Orden de visualización \\
        created\_at & DATETIME & Fecha de creación \\
        \bottomrule
    \end{tabular}
\end{table}

% -----------------------------------------------------------------------------
% ORDER
% -----------------------------------------------------------------------------
\subsection{Order (Pedido)}
\label{sec:order}

\begin{table}[ht]
    \centering
    \caption{Tabla: Order}
    \begin{tabular}{@{}lll@{}}
        \toprule
        \textbf{Atributo} & \textbf{Tipo} & \textbf{Descripción} \\
        \midrule
        id & BIGINT (PK) & Identificador único \\
        company\_id & BIGINT (FK) & Empresa asociada \\
        user\_id & BIGINT (FK) & Usuario que creó el pedido \\
        total\_amount & DECIMAL(10,2) & Monto total \\
        status & ENUM & Estado del pedido \\
        shipping\_address & VARCHAR(500) & Dirección de envío \\
        notes & TEXT & Notas adicionales \\
        created\_at & DATETIME & Fecha de creación \\
        updated\_at & DATETIME & Fecha de actualización \\
        \bottomrule
    \end{tabular}
\end{table}

\textbf{Estados del pedido:}
\begin{itemize}
    \item CART - En carrito
    \item PENDING\_PAYMENT - Pendiente de pago
    \item PAID - Pagado
    \item SHIPPED - Enviado
    \item DELIVERED - Entregado
    \item CANCELLED - Cancelado
\end{itemize}

\textbf{Relaciones:}
\begin{itemize}
    \item 1:N con OrderItem.
    \item N:1 con Company.
    \item N:1 con User.
    \item 1:1 con Payment.
\end{itemize}

% -----------------------------------------------------------------------------
% ORDERITEM
% -----------------------------------------------------------------------------
\subsection{OrderItem (Ítem del Pedido)}
\label{sec:orderitem}

\begin{table}[ht]
    \centering
    \caption{Tabla: OrderItem}
    \begin{tabular}{@{}lll@{}}
        \toprule
        \textbf{Atributo} & \textbf{Tipo} & \textbf{Descripción} \\
        \midrule
        id & BIGINT (PK) & Identificador único \\
        order\_id & BIGINT (FK) & Pedido asociado \\
        product\_id & BIGINT (FK) & Producto asociado \\
        product\_name & VARCHAR(255) & Snapshot nombre \\
        product\_sku & VARCHAR(50) & Snapshot SKU \\
        quantity & INT & Cantidad \\
        unit\_price & DECIMAL(10,2) & Precio unitario (snapshot) \\
        subtotal & DECIMAL(10,2) & Subtotal \\
        \bottomrule
    \end{tabular}
\end{table}

% -----------------------------------------------------------------------------
% PAYMENT
% -----------------------------------------------------------------------------
\subsection{Payment (Pago)}
\label{sec:payment}

\begin{table}[ht]
    \centering
    \caption{Tabla: Payment}
    \begin{tabular}{@{}lll@{}}
        \toprule
        \textbf{Atributo} & \textbf{Tipo} & \textbf{Descripción} \\
        \midrule
        id & BIGINT (PK) & Identificador único \\
        order\_id & BIGINT (FK) & Pedido asociado \\
        amount & DECIMAL(10,2) & Monto del pago \\
        provider & VARCHAR(50) & Proveedor (STRIPE, MERCADOPAGO) \\
        provider\_ref & VARCHAR(255) & Referencia del proveedor \\
        payment\_method & VARCHAR(50) & Método de pago \\
        status & ENUM & Estado del pago \\
        created\_at & DATETIME & Fecha de creación \\
        paid\_at & DATETIME & Fecha de pago \\
        \bottomrule
    \end{tabular}
\end{table}

\textbf{Estados del pago:}
\begin{itemize}
    \item PENDING - Pendiente
    \item PAID - Pagado
    \item FAILED - Fallido
    \item REFUNDED - Reembolsado
\end{itemize}

% -----------------------------------------------------------------------------
% MEDIA
% -----------------------------------------------------------------------------
\subsection{Media (Archivos)}
\label{sec:media}

\begin{table}[ht]
    \centering
    \caption{Tabla: Media}
    \begin{tabular}{@{}lll@{}}
        \toprule
        \textbf{Atributo} & \textbf{Tipo} & \textbf{Descripción} \\
        \midrule
        id & BIGINT (PK) & Identificador único \\
        url & VARCHAR(500) & URL del archivo \\
        filename & VARCHAR(255) & Nombre del archivo \\
        type & VARCHAR(50) & Tipo MIME \\
        size & BIGINT & Tamaño en bytes \\
        user\_id & BIGINT (FK) & Usuario que subió \\
        created\_at & DATETIME & Fecha de creación \\
        \bottomrule
    \end{tabular}
\end{table}

% -----------------------------------------------------------------------------
% REFRESHTOKEN
% -----------------------------------------------------------------------------
\subsection{RefreshToken (Token de Renovación)}
\label{sec:refreshtoken}

\begin{table}[ht]
    \centering
    \caption{Tabla: RefreshToken}
    \begin{tabular}{@{}lll@{}}
        \toprule
        \textbf{Atributo} & \textbf{Tipo} & \textbf{Descripción} \\
        \midrule
        id & BIGINT (PK) & Identificador único \\
        user\_id & BIGINT (FK) & Usuario asociado \\
        token & VARCHAR(500) & Token de renovación \\
        expires\_at & DATETIME & Fecha de expiración \\
        created\_at & DATETIME & Fecha de creación \\
        revoked & BOOLEAN & Estado revocado \\
        \bottomrule
    \end{tabular}
\end{table}

% -----------------------------------------------------------------------------
% AUDITLOG
% -----------------------------------------------------------------------------
\subsection{AuditLog (Registro de Auditoría)}
\label{sec:auditlog}

\begin{table}[ht]
    \centering
    \caption{Tabla: AuditLog}
    \begin{tabular}{@{}lll@{}}
        \toprule
        \textbf{Atributo} & \textbf{Tipo} & \textbf{Descripción} \\
        \midrule
        id & BIGINT (PK) & Identificador único \\
        user\_id & BIGINT (FK) & Usuario que realizó la acción \\
        action & VARCHAR(50) & Acción realizada \\
        entity & VARCHAR(50) & Entidad afectada \\
        entity\_id & BIGINT & ID de la entidad \\
        old\_value & JSON & Valor anterior \\
        new\_value & JSON & Valor nuevo \\
        ip\_address & VARCHAR(45) & Dirección IP \\
        user\_agent & VARCHAR(500) & Navegador/dispositivo \\
        timestamp & DATETIME & Fecha y hora \\
        \bottomrule
    \end{tabular}
\end{table}

% -----------------------------------------------------------------------------
% CONTENTBLOCK
% -----------------------------------------------------------------------------
\subsection{ContentBlock (Bloque de Contenido - CMS)}
\label{sec:contentblock}

\begin{table}[ht]
    \centering
    \caption{Tabla: ContentBlock}
    \begin{tabular}{@{}lll@{}}
        \toprule
        \textbf{Atributo} & \textbf{Tipo} & \textbf{Descripción} \\
        \midrule
        id & BIGINT (PK) & Identificador único \\
        section & VARCHAR(50) & Sección (homepage, about, etc.) \\
        key & VARCHAR(100) & Clave única \\
        content & TEXT & Contenido de texto \\
        image\_url & VARCHAR(500) & URL de imagen \\
        video\_url & VARCHAR(500) & URL de video \\
        sort\_order & INT & Orden de visualización \\
        active & BOOLEAN & Estado activo/inactivo \\
        created\_at & DATETIME & Fecha de creación \\
        updated\_at & DATETIME & Fecha de actualización \\
        \bottomrule
    \end{tabular}
\end{table}

% -----------------------------------------------------------------------------
% SETTING
% -----------------------------------------------------------------------------
\subsection{Setting (Configuración)}
\label{sec:setting}

\begin{table}[ht]
    \centering
    \caption{Tabla: Setting}
    \begin{tabular}{@{}lll@{}}
        \toprule
        \textbf{Atributo} & \textbf{Tipo} & \textbf{Descripción} \\
        \midrule
        id & BIGINT (PK) & Identificador único \\
        key & VARCHAR(100) & Clave de configuración \\
        value & TEXT & Valor (texto o JSON) \\
        type & ENUM & Tipo (STRING, NUMBER, BOOLEAN, JSON) \\
        description & VARCHAR(255) & Descripción \\
        created\_at & DATETIME & Fecha de creación \\
        updated\_at & DATETIME & Fecha de actualización \\
        \bottomrule
    \end{tabular}
\end{table}

% -----------------------------------------------------------------------------
% NOTIFICATIONTEMPLATE
% -----------------------------------------------------------------------------
\subsection{NotificationTemplate (Plantilla de Notificación)}
\label{sec:notificationtemplate}

\begin{table}[ht]
    \centering
    \caption{Tabla: NotificationTemplate}
    \begin{tabular}{@{}lll@{}}
        \toprule
        \textbf{Atributo} & \textbf{Tipo} & \textbf{Descripción} \\
        \midrule
        id & BIGINT (PK) & Identificador único \\
        type & VARCHAR(50) & Tipo de notificación \\
        subject & VARCHAR(255) & Asunto \\
        body & TEXT & Cuerpo (con placeholders) \\
        is\_active & BOOLEAN & Esta activo \\
        created\_at & DATETIME & Fecha de creación \\
        updated\_at & DATETIME & Fecha de actualización \\
        \bottomrule
    \end{tabular}
\end{table}

% -----------------------------------------------------------------------------
% USERNOTIFICATION
% -----------------------------------------------------------------------------
\subsection{UserNotification (Notificación de Usuario)}
\label{sec:usernotification}

\begin{table}[ht]
    \centering
    \caption{Tabla: UserNotification}
    \begin{tabular}{@{}lll@{}}
        \toprule
        \textbf{Atributo} & \textbf{Tipo} & \textbf{Descripción} \\
        \midrule
        id & BIGINT (PK) & Identificador único \\
        user\_id & BIGINT (FK) & Usuario destinatario \\
        title & VARCHAR(255) & Título \\
        message & TEXT & Mensaje \\
        type & ENUM & Tipo (INFO, SUCCESS, WARNING, ERROR) \\
        is\_read & BOOLEAN & Leída \\
        link\_url & VARCHAR(500) & Enlace relacionado \\
        created\_at & DATETIME & Fecha de creación \\
        \bottomrule
    \end{tabular}
\end{table}

% -----------------------------------------------------------------------------
% COMPANYADDRESS
% -----------------------------------------------------------------------------
\subsection{CompanyAddress (Dirección de Empresa)}
\label{sec:companyaddress}

\begin{table}[ht]
    \centering
    \caption{Tabla: CompanyAddress}
    \begin{tabular}{@{}lll@{}}
        \toprule
        \textbf{Atributo} & \textbf{Tipo} & \textbf{Descripción} \\
        \midrule
        id & BIGINT (PK) & Identificador único \\
        company\_id & BIGINT (FK) & Empresa asociada \\
        label & VARCHAR(50) & Etiqueta (Principal, Bodega, etc.) \\
        address & VARCHAR(500) & Dirección \\
        city & VARCHAR(100) & Ciudad \\
        department & VARCHAR(100) & Departamento \\
        zip\_code & VARCHAR(20) & Código postal \\
        is\_default & BOOLEAN & Es dirección predeterminada \\
        active & BOOLEAN & Estado activo \\
        created\_at & DATETIME & Fecha de creación \\
        \bottomrule
    \end{tabular}
\end{table}

% -----------------------------------------------------------------------------
% SHIPPINGZONE
% -----------------------------------------------------------------------------
\subsection{ShippingZone (Zona de Envío)}
\label{sec:shippingzone}

\begin{table}[ht]
    \centering
    \caption{Tabla: ShippingZone}
    \begin{tabular}{@{}lll@{}}
        \toprule
        \textbf{Atributo} & \textbf{Tipo} & \textbf{Descripción} \\
        \midrule
        id & BIGINT (PK) & Identificador único \\
        name & VARCHAR(100) & Nombre de la zona \\
        department & VARCHAR(100) & Departamento/Región \\
        delivery\_cost & DECIMAL(10,2) & Costo de envío \\
        estimated\_days & INT & Días estimados \\
        active & BOOLEAN & Estado activo/inactivo \\
        \bottomrule
    \end{tabular}
\end{table}

% -----------------------------------------------------------------------------
% FAVORITEPRODUCT
% -----------------------------------------------------------------------------
\subsection{FavoriteProduct (Producto Favorito)}
\label{sec:favoriteproduct}

\begin{table}[ht]
    \centering
    \caption{Tabla: FavoriteProduct}
    \begin{tabular}{@{}lll@{}}
        \toprule
        \textbf{Atributo} & \textbf{Tipo} & \textbf{Descripción} \\
        \midrule
        id & BIGINT (PK) & Identificador único \\
        user\_id & BIGINT (FK) & Usuario \\
        product\_id & BIGINT (FK) & Producto \\
        created\_at & DATETIME & Fecha de creación \\
        \bottomrule
    \end{tabular}
\end{table}

% -----------------------------------------------------------------------------
% QUOTE
% -----------------------------------------------------------------------------
\subsection{Quote (Cotización)}
\label{sec:quote}

\begin{table}[ht]
    \centering
    \caption{Tabla: Quote}
    \begin{tabular}{@{}lll@{}}
        \toprule
        \textbf{Atributo} & \textbf{Tipo} & \textbf{Descripción} \\
        \midrule
        id & BIGINT (PK) & Identificador único \\
        company\_id & BIGINT (FK) & Empresa solicitante \\
        user\_id & BIGINT (FK) & Usuario solicitante \\
        status & ENUM & Estado (PENDING, APPROVED, REJECTED, EXPIRED) \\
        notes & TEXT & Notas adicionales \\
        valid\_until & DATE & Validez hasta \\
        created\_at & DATETIME & Fecha de creación \\
        updated\_at & DATETIME & Fecha de actualización \\
        \bottomrule
    \end{tabular}
\end{table}

% -----------------------------------------------------------------------------
% QUOTEITEM
% -----------------------------------------------------------------------------
\subsection{QuoteItem (Ítem de Cotización)}
\label{sec:quoteitem}

\begin{table}[ht]
    \centering
    \caption{Tabla: QuoteItem}
    \begin{tabular}{@{}lll@{}}
        \toprule
        \textbf{Atributo} & \textbf{Tipo} & \textbf{Descripción} \\
        \midrule
        id & BIGINT (PK) & Identificador único \\
        quote\_id & BIGINT (FK) & Cotización asociada \\
        product\_id & BIGINT (FK) & Producto cotizado \\
        quantity & INT & Cantidad solicitada \\
        unit\_price & DECIMAL(10,2) & Precio cotizado \\
        notes & TEXT & Notas \\
        \bottomrule
    \end{tabular}
\end{table}

% -----------------------------------------------------------------------------
% SALESREPORT
% -----------------------------------------------------------------------------
\subsection{SalesReport (Reporte de Ventas)}
\label{sec:salesreport}

\begin{table}[ht]
    \centering
    \caption{Tabla: SalesReport}
    \begin{tabular}{@{}lll@{}}
        \toprule
        \textbf{Atributo} & \textbf{Tipo} & \textbf{Descripción} \\
        \midrule
        id & BIGINT (PK) & Identificador único \\
        report\_type & ENUM & Tipo (DAILY, MONTHLY, YEARLY, CUSTOM) \\
        company\_id & BIGINT (FK) & Empresa (null para todos) \\
        start\_date & DATE & Fecha de inicio \\
        end\_date & DATE & Fecha de fin \\
        total\_orders & INT & Total de pedidos \\
        total\_revenue & DECIMAL(15,2) & Ingresos totales \\
        total\_products\_sold & INT & Productos vendidos \\
        average\_order\_value & DECIMAL(10,2) & Valor promedio \\
        top\_products & JSON & Productos más vendidos \\
        generated\_by & BIGINT (FK) & Usuario que generó \\
        created\_at & DATETIME & Fecha de creación \\
        \bottomrule
    \end{tabular}
\end{table}

% ==============================================================================
% CAPÍTULO 6: FLUJOS DEL SISTEMA
% ==============================================================================
\chapter{Flujos del Sistema}
\label{cap:flujos}

\section{Flujo de Autenticación}
\label{flujo:auth}

\begin{enumerate}
    \item Usuario se registra o inicia sesión.
    \item \textbf{(Opcional) 2FA}: Si está habilitado, solicitar código TOTP.
    \item Generar access token + refresh token.
    \item Acceso al sistema.
\end{enumerate}

\section{Flujo de Recuperación de Contraseña}
\label{flujo:password}

\begin{enumerate}
    \item Usuario solicita ``Olvidé mi contraseña''.
    \item Sistema genera token de recuperación (expiración: 1 hora).
    \item Email con enlace de recuperación.
    \item Usuario ingresa nueva contraseña.
    \item Invalidar tokens anteriores.
\end{enumerate}

\section{Flujo de Compra}
\label{flujo:compra}

\begin{enumerate}
    \item Empresa se registra en el sistema.
    \item Usuario inicia sesión con JWT.
    \item Explora productos (catálogo dinámico desde BD).
    \item Agrega productos al carrito.
    \item Crea pedido.
    \item Paga por pasarela.
    \item Pedido confirmado.
    \item Auditoría registrada.
\end{enumerate}

\section{Flujo 2FA (Autenticación de Dos Factores)}
\label{flujo:2fa}

\begin{enumerate}
    \item Usuario habilita 2FA desde su perfil.
    \item Sistema genera secreto TOTP.
    \item Se muestra QR para Google Authenticator.
    \item Usuario verifica con código inicial.
    \item En login, después de contraseña:
    \begin{itemize}
        \item Si 2FA habilitado $\rightarrow$ solicitar código.
        \item Si no $\rightarrow$ continuar normalmente.
    \end{itemize}
\end{enumerate}

\section{Flujo Cambio de Contraseña}
\label{flujo:cambio_password}

\begin{enumerate}
    \item Usuario va a su perfil.
    \item Clic en ``Cambiar contraseña''.
    \item Ingresa contraseña actual + nueva contraseña.
    \item Sistema valida contraseña actual.
    \item Actualiza hash de contraseña.
    \item \textbf{Invalidar todos los tokens activos} (logout de todos los dispositivos).
    \item Enviar email de notificación.
\end{enumerate}

\section{Flujo Cerrar Sesión (Logout)}
\label{flujo:logout}

\begin{enumerate}
    \item Usuario hace clic en ``Cerrar sesión''.
    \item Frontend envía solicitud al backend.
    \item Backend marca refresh token como revoke.
    \item Frontend elimina access token del storage.
    \item Redireccionar a login.
    \item (Opcional) Opción de ``Cerrar todas las sesiones''.
\end{enumerate}

% ==============================================================================
% CAPÍTULO 7: PANEL DE ADMINISTRACIÓN
% ==============================================================================
\chapter{Panel de Administración}
\label{cap:admin}

\section{Funcionalidades del Administrador}

\subsection{Gestión de Contenido (CMS)}
\begin{itemize}
    \item Editar texto de página principal.
    \item Gestionar banners e imágenes.
    \item Actualizar About Us, Contacto.
    \item Editar políticas y términos.
    \item Gestionar footer.
\end{itemize}

\subsection{Gestión de Productos}
\begin{itemize}
    \item CRUD completo de productos.
    \item CRUD de categorías.
    \item Subida de imágenes múltiples.
    \item Gestionar precio base y precios por volumen.
    \item Control de inventario.
\end{itemize}

\subsection{Gestión de Empresas y Usuarios}
\begin{itemize}
    \item Ver empresas registradas.
    \item Activar/desactivar empresas.
    \item Gestionar usuarios por empresa.
    \item Asignar roles.
\end{itemize}

\subsection{Gestión de Pedidos}
\begin{itemize}
    \item Ver todos los pedidos.
    \item Cambiar estados.
    \item Procesar devoluciones.
\end{itemize}

\subsection{Configuración de Seguridad}
\begin{itemize}
    \item Habilitar/deshabilitar 2FA.
    \item Configurar políticas de contraseña.
    \item Ver logs de auditoría.
\end{itemize}

\subsection{Dashboard y Reportes}
\begin{itemize}
    \item \textbf{Estadísticas de ventas}
    \begin{itemize}
        \item Gráfico de ventas diarias/semanales/mensuales.
        \item Ingresos totales.
        \item Cantidad de pedidos.
        \item Valor promedio por pedido.
    \end{itemize}
    \item \textbf{Reportes por empresa}
    \begin{itemize}
        \item Ventas por empresa.
        \item Empresas más activas.
        \item Empresas nuevas.
    \end{itemize}
    \item \textbf{Reporte de inventario}
    \begin{itemize}
        \item Productos más vendidos.
        \item Stock bajo.
        \item Productos sin movimiento.
    \end{itemize}
    \item \textbf{KPI's principales}
    \begin{itemize}
        \item Ticket promedio.
        \item Tasa de conversión.
        \item Usuarios activos.
    \end{itemize}
    \item \textbf{Exportar reportes} (Excel, PDF).
\end{itemize}

\subsection{Gestión de Cotizaciones}
\begin{itemize}
    \item Ver solicitudes de cotización.
    \item Aprobar/rechazar cotizaciones.
    \item Definir precios especiales.
    \item Enviar cotizaciones por email.
\end{itemize}

\subsection{Notificaciones del Sistema}
\begin{itemize}
    \item Enviar notificaciones a usuarios.
    \item Notificaciones masivas.
    \item Historial de notificaciones.
\end{itemize}

% ==============================================================================
% CAPÍTULO 8: NOTAS DE IMPLEMENTACIÓN
% ==============================================================================
\chapter{Notas de Implementación}
\label{cap:implementacion}

\section{Seguridad}
\begin{itemize}
    \item Usar \textbf{bcrypt} o \textbf{Argon2id} para contraseñas.
    \item Tokens JWT cortos (15-30 min access, 7 días refresh).
    \item Almacenar refresh tokens en BD con capacidad de revocación.
    \item Registrar todas las acciones sensibles.
\end{itemize}

\section{CMS}
\begin{itemize}
    \item ContentBlock usa pares section/key para identificar contenido.
    \item Frontend consume estos valores dinámicamente.
    \item Admin panel proporciona editor WYSIWYG simple.
\end{itemize}

\section{Productos}
\begin{itemize}
    \item Precio base + precios por cantidad (price\_tiers JSON).
    \item SKU único por producto.
    \item Control de stock transaccional.
\end{itemize}

% ==============================================================================
% CAPÍTULO 9: ALCANCE MVP (1 MES)
% ==============================================================================
\chapter{Alcance MVP - Versión 1.0}
\label{cap:mvp}

\alert{IMPORTANTE:} 1 mes solo = Proyecto reducido obligatorio

\section{Semana 1: Base y Autenticación}
\begin{itemize}
    \item[ ] Estructura del proyecto Spring Boot + React.
    \item[ ] Entities básicas: User, Company, Role.
    \item[ ] Login JWT sin 2FA.
    \item[ ] Registro de empresas.
    \item[ ] Logout básico.
\end{itemize}

\section{Semana 2: Productos y Carrito}
\begin{itemize}
    \item[ ] CRUD Products + Categories.
    \item[ ] Imágenes (solo 1 por producto).
    \item[ ] Carrito en memoria (no persistente).
    \item[ ] Crear Order.
\end{itemize}

\section{Semana 3: Pagos e Inventario}
\begin{itemize}
    \item[ ] Integrar 1 pasarela (Stripe o MercadoPago).
    \item[ ] Webhook de confirmación.
    \item[ ] Actualizar stock al comprar.
    \item[ ] Historial de pedidos.
\end{itemize}

\section{Semana 4: Admin Básico y Pulir}
\begin{itemize}
    \item[ ] Admin: ver pedidos.
    \item[ ] Admin: activar/desactivar productos.
    \item[ ] Login empresa básico.
    \item[ ] Deploy básico.
\end{itemize}

\section{No Incluido en V1}
\begin{alert}{FUERA DEL ALCANCE MVP:}
\begin{itemize}
    \item 2FA.
    \item CMS/ContentBlock.
    \item Dashboard de reportes.
    \item Cotizaciones B2B.
    \item Wishlist.
    \item Carrito persistente.
    \item Múltiples imágenes.
    \item Notificaciones.
    \item Configuraciones globales.
\end{itemize}
\end{alert}

\section{Prioridades de Desarrollo}
\begin{enumerate}
    \item \textbf{Día 1-7}: Login + Registro Empresa + Productos.
    \item \textbf{Día 8-14}: Carrito + Order + Payment.
    \item \textbf{Día 15-21}: Admin básico + Inventario.
    \item \textbf{Día 22-30}: Testing + Fixes + Deploy.
\end{enumerate}

% ==============================================================================
% APÉNDICES
% ==============================================================================
\appendix
\chapter{Tablas Requeridas del Sistema}
\label{apendice:tablas}

\begin{table}[ht]
    \centering
    \caption{Lista completa de tablas del sistema}
    \begin{tabular}{cll}
        \toprule
        \textbf{Nº} & \textbf{Tabla} & \textbf{Descripción} \\
        \midrule
        1 & Company & Empresas registradas \\
        2 & User & Usuarios del sistema (con campos 2FA) \\
        3 & Role & Roles del sistema \\
        4 & UserRole & Relación usuario-rol \\
        5 & Category & Categorías de productos \\
        6 & Product & Productos disponibles \\
        7 & ProductImage & Imágenes de productos \\
        8 & Order & Pedidos de compra \\
        9 & OrderItem & Ítems del pedido \\
        10 & Payment & Registro de pagos \\
        11 & Media & Archivos multimedia \\
        12 & RefreshToken & Tokens de renovación \\
        13 & AuditLog & Registro de auditoría \\
        14 & ContentBlock & Contenido CMS \\
        15 & Setting & Configuraciones globales \\
        16 & NotificationTemplate & Plantillas de notificación \\
        17 & UserNotification & Notificaciones de usuario \\
        18 & CompanyAddress & Direcciones de empresa \\
        19 & ShippingZone & Zonas de envío \\
        20 & FavoriteProduct & Productos favoritos \\
        21 & Quote & Cotizaciones B2B \\
        22 & QuoteItem & Ítems de cotización \\
        23 & SalesReport & Reportes de ventas \\
        \bottomrule
    \end{tabular}
\end{table}

% ==============================================================================
% REFERENCIAS
% ==============================================================================
\chapter{Referencias}
\label{cap:referencias}

\begin{thebibliography}{99}
    
    \bibitem{spring-boot} 
    \textbf{Spring Boot Reference Documentation}\\
    \textit{Pivotal Software, Inc.}\\
    \href{https://spring.io/projects/spring-boot}{https://spring.io/projects/spring-boot}
    
    \bibitem{spring-security}
    \textbf{Spring Security Reference}\\
    \textit{Pivotal Software, Inc.}\\
    \href{https://spring.io/projects/spring-security}{https://spring.io/projects/spring-security}
    
    \bibitem{jwt}
    \textbf{JWT (JSON Web Tokens)}\\
    \textit{Internet Engineering Task Force (IETF)}\\
    \href{https://jwt.io}{https://jwt.io}
    
    \bibitem{totp}
    \textbf{TOTP: Time-based One-Time Passwords}\\
    \textit{RFC 6238}\\
    \href{https://tools.ietf.org/html/rfc6238}{https://tools.ietf.org/html/rfc6238}
    
    \bibitem{react}
    \textbf{React Documentation}\\
    \textit{Meta Platforms, Inc.}\\
    \href{https://react.dev}{https://react.dev}
    
    \bibitem{vite}
    \textbf{Vite - Next Generation Frontend Tooling}\\
    \textit{Evan You}\\
    \href{https://vitejs.dev}{https://vitejs.dev}
    
    \bibitem{mercado-pago}
    \textbf{MercadoPago Developers}\\
    \textit{Mercado Libre}\\
    \href{https://www.mercadopago.com.co/developers}{https://www.mercadopago.com.co/developers}
    
    \bibitem{bcrypt}
    \textbf{BCrypt Password Hashing Algorithm}\\
    \textit{ProPublica}\\
    \href{https://www.propublica.org/article/bcrypt-password-hashing}{https://www.propublica.org}
    
\end{thebibliography}

% ==============================================================================
% FIN DEL DOCUMENTO
% ==============================================================================
\end{document}
